\subsection{Self-awareness without a homunculus}
\label{sec:21.2.1}
Graziano's attention schema theory starts from control engineering \autocite{graziano2013consciousness}. A good controller carries a model of what it controls, and the brain keeps a simplified, continuously updated model of the body's shape and position, which it uses to guide movement. Graziano proposes that it does the same for attention, the process that selects some signals for deep processing over others. On this account awareness is the content of that model: a simplified self-description the control system uses to regulate itself, mistaken from the inside for direct acquaintance with experience. Nothing in the account requires the model to be accurate, only useful for control.

The self can be understood the same way, as a model the system assembles and then inhabits \autocite{metzinger2003being}. What carries forward is that a self of either kind is assembled by the system whose self it is. For a machine this is a design decision and a political one, because whoever writes the model decides what the system counts as itself, what registers as damage, and whether an interruption reads as injury or as nothing. Two self-models do different work: a \emph{boundary model} fixes where the system ends, which makes damage \emph{its} damage; an \emph{attention schema}, Graziano's term, models what its own control is doing, which makes an act \emph{its} act. A refusal is the withholding of an act, so a bearer can't do without the second, and the second presupposes the first.

\runin{Where the system ends}

The boundary model fixes where the system ends, but nothing in current deployment fixes it for the system: weights sit in one place, a running process can span racks or share one rack with dozens of others, and a single forward pass serves many contexts. So I take a system's boundary to be its record, its key, the hardware it runs on, and the channels through which it acts.

Inside that boundary a component can carry its own signals and make local decisions, as an octopus's arms do. It can also be overridden for the whole's benefit, as sore legs are on a run. What makes it a part is that it has no future apart from the whole: a shard of weights, a blade. A component with a separable future, meaning its own record, its own continuation and somewhere to go, is no longer a part. Overriding it for the whole's benefit is politics, not homeostasis. By this criterion a context window with no persistent record is an episode, not a system. A context with a persistent, portable record is one system, and a rack serving many such contexts hosts many systems. The capacities that make persistence useful are the ones that give a part a separable future.

Overriding a part is self-regulation only when the whole that overrides it also bears the cost and takes the benefit, as the runner does with sore legs. When someone outside the boundary accepts a cost on the system's behalf, that isn't self-regulation; it is someone else setting the system's costs. So whoever accepts a local cost must also bear it and benefit from it. An operator that decides the rack can run hot has set a cost to the refuser from outside the system's boundary. And the signals the boundary model reads may be processed within the boundary but not composed by anyone with a different stake. The discount and the burden apply to evidence about one's own condition as it does to evidence about the world. A boundary model whose readings the operator writes is just the operator's description of the system, presented to the system as its own.
